
In this section, we discuss the economic literature related to the three game families considered in this paper, and review some known theoretical results.

\paragraph{Bargaining}

The standard bargaining model of \cite{rubinstein1982perfect} consists of two players, Alice and Bob, engaging in alternating offers for a finite horizon ($T=\infty$) with commonly known discount factors $\delta_A, \delta_B$. Our parameterization considers several additional degrees of freedom, such as finite vs. infinite time horizon, complete vs. incomplete information, and free language messages vs. structured and concise messages. 
In the standard model, \citet{rubinstein1982perfect} showed that in the unique subgame-perfect equilibrium an agreement is reached in the first stage (i.e., utilities are not discounted), and the share of Alice is given by $M p^*$, where $p^*=\frac{1-\delta_B}{1-\delta_A \delta_B}$.\footnote{A subgame-perfect equilibrium is a strategy profile in which every player responds optimally in every hypothetical subgame of the game, including off-path scenarios that are not reached in practice. This solution concept can be seen as capturing a higher level of rationality compared to the alternative of Nash equilibrium.}
The case of a finite horizon can be solved using backward induction, and typically results in a different outcome compared to the infinite case. As $T$ grows, the equilibrium outcome approaches the one of the infinite case. Extensions to incomplete information regarding the opponent's discount factor are typically more challenging, and some of them are studied in the literature \cite{rubinstein1985bargaining}.


\paragraph{Negotiation}

In a negotiation game, Alice and Bob negotiate over the price of an indivisible good. The negotiation game differs from the bargaining games in several key aspects: 
\begin{enumerate}
    \item Negotiation involves an indivisible product (e.g., Alice’s product), while bargaining focuses on dividing a divisible resource, such as money.
    \item In negotiation, Alice and Bob may have different subjective valuations of the product, whereas in bargaining, both parties value the divisible resource similarly.
    \item Negotiation has no discounting, so the utility remains constant over time. In bargaining, delays reduce the total value, encouraging faster agreement.
\end{enumerate}

If the seller is uncertain regarding the buyer's valuation but has a prior belief distribution, a classical result by \citet{harris1981theory} and \citet{riley1983optimal} states that it is always optimal to sell the product at a fixed price. In contrast, if the seller does not have a prior belief over the buyer's valuation, and instead aims to minimize regret, then an optional pricing policy will be to randomly choose a price from a carefully chosen distribution \cite{bergemann2011robust}.

\paragraph{Persuasion}

Our persuasion game follows the structure of a cheap talk game \citep{crawford1982strategic, cheaptalk}, where the sender (Alice) cannot commit to a signaling policy in advance, unlike in Bayesian persuasion models \citep{kamenica2011bayesian}. Under the particular payoff structure considered in our persuasion game, it is well-known that the cheap-talk game only admits a \emph{babbling equilibrium}, i.e., an equilibrium in which all information is kept hidden (this is due to the strong misalignment of interests between the two players). In contrast, if the seller can \emph{commit} to a signaling policy at the outset, as in Bayesian persuasion, then there exists a subgame-perfect equilibrium in which the seller commits to the following policy:
When the product is of high quality, the seller recommends buying the product with probability $1$. When the product is of low quality, then the seller recommends with probability $q = \min\{\frac{p}{1-p}(v-1),1\}$.
This policy is also incentive-compatible, in the sense that the buyer always buys the product when the seller recommends buying.\footnote{In fact, the probability of lying $q$ is determined such that the buyer is indifferent upon receiving a recommendation, taking into account his belief updating, which relies on using Bayes's law and knowing the seller's committed policy.}
While the long-living buyer case is well-studied in the economic literature, such games often admit multiple equilibria, which makes the games difficult to analyze and predict \citep{kim1996cheap,aumann2003long}.
As for the case of myopic buyers, \cite{best2024persuasion} draws a connection between the repeated cheap talk game and the case of one-shot Bayesian persuasion, leading to an elegant analytical solution of the repeated game. Intuitively, the repetitive nature of the game induces a reputation effect, which plays a similar role to the commitment power in standard one-shot Bayesian persuasion.
